\section{Problem setup and cost accounting}
\label{sec:setup}

\subsection{Newton--CG step and action space}

At each step we solve the damped Newton system with truncated conjugate
gradients \citep{hestenes1952cg,steihaug1983cg}, forming no matrix and using
only Hessian-vector products \citep{pearlmutter1994hvp}:
\begin{equation}
  (H + \lambda I)\, p = -g,
  \qquad \text{terminated after $k$ CG iterations.}
\end{equation}

The controller chooses along two axes: a damping multiplier $m$, applied as
$\lambda \leftarrow \mathrm{clip}(\lambda \cdot m)$ and accumulated in log space,
and a CG iteration budget $k \in \{3, 5, 10, 20\}$.

The step size is fixed at $1.0$. It aliases with the damping axis, so leaving it
free would make the two axes' contributions inseparable.

We use three action spaces:

\begin{center}
\begin{tabular}{llr}
\toprule
Name & Actions & Damping range ($\log_{10}$) \\
\midrule
\ctrl{narrow}   & 12  & 0.95 \\
\ctrl{wide}     & 28  & 2.86 \\
\ctrl{absolute} & 132 & 15.27 \\
\bottomrule
\end{tabular}
\end{center}

\ctrl{narrow} is the deployment target; \ctrl{absolute} exists to remove
reachability constraints for analysis. \ctrl{absolute} has \emph{the same
logarithmic resolution} as \ctrl{narrow}, so that the comparison ``is the action
space the bottleneck?'' is not confounded by a difference in resolution.

\subsection{Cost accounting and its scope}
\label{sec:ge}

We do not use wall-clock time as the primary metric. At this scale (thousands to
tens of thousands of parameters) time is dominated by kernel launch and
interpreter overhead rather than floating-point work, so measurements would
reflect implementation overhead rather than optimization efficiency.

We use a hardware-independent unit instead. One gradient-equivalent (GE) is one
forward plus one backward pass over a gradient batch, and the per-step cost is
\begin{equation}
  \mathrm{cost}_{\GE}(k) = c_{\text{grad-graph}} + k \cdot c_{\text{hvp}}
    + c_{\text{fwd}} .
\end{equation}

All comparisons are made at an \emph{identical GE budget}. In Newton--CG the
per-step cost varies strongly with the action ($k = 3$ versus $k = 20$ differs by
more than a factor of six), so matching step counts and comparing final loss
would compare runs of unequal cost.

We state the scope of this definition in advance.

\begin{quote}
GE normalizes gradient-equivalent oracle calls within a regime, not total
floating-point operations across different batch sizes. Cross-regime absolute
performance comparisons are therefore descriptive rather than compute-normalized.
\end{quote}

The quadratic comparisons of Section~\ref{sec:heldout} all take place on the same
objective and are unaffected by this limitation. It applies only to absolute
cross-regime comparisons in Section~\ref{sec:mismatch}.

\subsection{Execution environment and the nature of GE}
\label{sec:cpu}

\begin{quote}
All Stage~2 optimization and planning experiments were executed on CPU. The
GPU-derived cost-model files in the repository originate from an earlier Stage~1
calibration and were not used in Stage~2; Stage~2 GE accounting used
HVP-equivalent oracle counts.
\end{quote}

The repository retains those Stage~1 configuration files, including their
\texttt{device: cuda:0} entries, for provenance. They do not describe the Stage~2
execution environment.

Reading the search cost also requires a distinction.

\begin{quote}
Search GE measures simulated oracle work rather than wall-clock GPU compute.
\end{quote}

The factor of $1{,}294$ reported in Section~\ref{sec:cost} is a ratio of
\emph{oracle calls} performed by the planner, not a ratio of GPU wall-clock time.

\subsection{Evaluation metric}

We measure improvement at a fixed budget as a logarithmic improvement in nats,
\begin{equation}
  \JE = \log L_0 - \log L_{\text{final}} .
\end{equation}

We place a numerical floor at a relative value of $100 \cdot \varepsilon$ and cap
values below it. We report both the uncapped pairs and the full set of pairs.
Whenever a pair is dropped from a statistic, the reason is recorded.
